%%%%%%%%%%%%%%%%%%%%%%%%%%%%%%%%%%%%%%%%%%%%%%%%%%%%%%%%%%%%%%%%%%%%%%%%%%%%%%%%
%2345678901234567890123456789012345678901234567890123456789012345678901234567890
%        1         2         3         4         5         6         7         8

\documentclass[letterpaper, 10 pt, conference]{ieeeconf}

\IEEEoverridecommandlockouts
\overrideIEEEmargins

% Required packages
\usepackage{amsmath}
\usepackage{amssymb}
\usepackage{bm}
\usepackage{booktabs}
\usepackage{graphicx}

\newtheorem{remark}{Remark}

\title{\LARGE \bf
Learning-Based Backward Reach-Avoid Tubes for Safe Spacecraft Proximity Docking in 13 Dimensions
}

\author{Albert Author$^{1}$ and Bernard D. Researcher$^{2}$%
\thanks{*This work was not supported by any organization}%
\thanks{$^{1}$Albert Author is with Faculty of Electrical Engineering, Mathematics and Computer Science,
        University of Twente, 7500 AE Enschede, The Netherlands
        {\tt\small albert.author@papercept.net}}%
\thanks{$^{2}$Bernard D. Researcheris with the Department of Electrical Engineering, Wright State University,
        Dayton, OH 45435, USA
        {\tt\small b.d.researcher@ieee.org}}%
}

\begin{document}

\maketitle
\thispagestyle{empty}
\pagestyle{empty}

%%%%%%%%%%%%%%%%%%%%%%%%%%%%%%%%%%%%%%%%%%%%%%%%%%%%%%%%%%%%%%%%%%%%%%%%%%%%%%%%
\begin{abstract}
Spacecraft docking demands control policies that jointly guarantee target reachability and collision avoidance under coupled translational--rotational dynamics, making it a natural application for Backward Reach-Avoid Tube (BRAT) analysis via Hamilton--Jacobi (HJ) methods. However, the full docking problem requires a 13-dimensional state space, far exceeding the reach of classical grid-based BRAT computation. In this paper, we develop a learning-based BRAT framework to bridge this gap. We approximate the HJ value function with a neural network, combining PDE-based training with curriculum-aware value labels generated via Model Predictive Control (MPC) in a semi-supervised scheme that improves accuracy where the PDE residual alone is insufficient. We introduce an exact-initial-condition value parameterization tailored to reach-avoid semantics. The resulting controller provides continuous safety guarantees over the planning horizon. We first validate the framework on a 6-dimensional planar docking model, recovering a solution similar to the grid-based reference solution, and then scale to the full 13-dimensional problem, a regime inaccessible to existing grid-based methods. Compared to MPC and Reinforcement Learning (RL) baselines, the proposed method achieves competitive or superior docking success at substantially lower online computational cost.
\end{abstract}

%%%%%%%%%%%%%%%%%%%%%%%%%%%%%%%%%%%%%%%%%%%%%%%%%%%%%%%%%%%%%%%%%%%%%%%%%%%%%%%%
\section{INTRODUCTION}

Autonomous spacecraft docking is a foundational capability for on-orbit servicing, debris removal, and space station resupply~\cite{fehse2003}. As missions grow in complexity---from NASA's Habitable Worlds Observatory requiring autonomous on-orbit assembly~\cite{hwo2024} to commercial satellite life-extension services---guaranteed safety during proximity operations has become a hard requirement rather than a desirable feature~\cite{ventura2017}. A docking maneuver must simultaneously steer the chaser spacecraft into a narrow capture corridor while avoiding collision with the target body, under coupled translational and rotational dynamics.

Hamilton--Jacobi (HJ) reachability analysis provides a principled framework for certifying such reach-avoid behaviors~\cite{mitchell2005, bansal2017hjsurvey}. The resulting value function encodes the set of all initial states from which the system can reach a goal while avoiding obstacles, and its gradient yields a feedback control law. However, classical grid-based HJ methods suffer from the curse of dimensionality: computational cost grows exponentially with state dimension, limiting practical application to systems of at most five to six dimensions~\cite{mitchell2005}. A full spacecraft docking model with three-dimensional translation, quaternion-based attitude representation, and their rates requires a 13-dimensional state space---far beyond any existing grid-based solver.

In this paper, we develop a learning-based Backward Reach-Avoid Tube (BRAT) framework that overcomes this dimensional barrier. We approximate the HJ value function with a neural network trained via a combination of PDE-based residual minimization and curriculum-aware supervision from MPC-generated value labels. A key ingredient is an exact-initial-condition value parameterization that enforces the terminal boundary condition by construction, eliminating a major source of training error for reach-avoid problems. The resulting value function provides a continuous safety metric: its sign indicates whether safe docking is achievable from a given state, and its gradient yields the optimal control action.

\subsection{Related Work}

\textit{HJ Reachability and Safety Certification:}
Classical HJ methods compute reachable sets by solving the Hamilton--Jacobi--Isaacs (HJI) PDE on a discrete grid~\cite{mitchell2005, lygeros2004}. The level-set framework established by Mitchell et al.~\cite{mitchell2005} underpins modern reachability tools such as \texttt{hj\_reachability}~\cite{hjreachability}. Recent work has applied these methods to spacecraft proximity operations: Bui et al.~\cite{bui2022} computed reach-avoid sets for four-dimensional space vehicle docking under continuous thrust, and Mote et al.~\cite{mote2023} addressed rendezvous with a tumbling target. However, grid-based approaches are fundamentally limited to low dimensions. Decomposition methods~\cite{chen2018decomposition} can handle higher dimensions when subsystems are decoupled, but the translational--rotational coupling in three-dimensional docking precludes direct decomposition of the full 13D system.

\textit{Learning-Based Reachability:}
DeepReach~\cite{bansal2021deepreach} demonstrated that neural networks with sinusoidal representation networks (SIREN)~\cite{sitzmann2020siren} can approximate HJ value functions in systems up to nine dimensions by minimizing the PDE residual as a self-supervised loss. Subsequent work has extended this framework to parameter-conditioned reachable sets~\cite{borquez2023} and safety-constrained trajectory design~\cite{shao2023}. However, existing DeepReach applications primarily address backward reachable tubes (BRTs) for avoid-only problems. The reach-avoid extension---requiring simultaneous goal reachability and obstacle avoidance---introduces a more complex variational inequality (VI) that has not been demonstrated at this scale.

\textit{Spacecraft Proximity Operations and Control:}
Traditional approaches to autonomous docking employ LQR controllers~\cite{jewison2016}, fuel-optimal convex programming~\cite{malyuta2022}, or MPC-based trajectory optimization~\cite{eren2015}. While computationally efficient online, these methods provide no formal characterization of the safe operating envelope. Reinforcement learning approaches~\cite{hovell2021, broida2019} have shown promising results for docking but similarly lack safety certificates. Safety filter architectures~\cite{wabersich2021, ames2019cbf} can overlay guarantees on nominal controllers but require a pre-computed safe set---precisely what our BRAT provides.

\subsection{Contributions}

Our contributions are as follows:
\begin{enumerate}
\item We formulate the full 13-dimensional spacecraft docking problem---comprising three-dimensional translation, quaternion attitude, and their rates---as a neural BRAT, including orientation-dependent collision geometry and attitude-aware goal and failure sets derived from International Docking System Standard (IDSS) tolerances~\cite{nasaidss}.
\item We develop a semi-supervised training pipeline that combines DeepReach PDE-based training with MPC-generated curriculum labels, employing an exact-initial-condition value parameterization tailored to reach-avoid semantics that eliminates the Dirichlet boundary loss entirely.
\item We validate the framework on a 6D planar docking model against grid-based ground truth, demonstrating recovery of the certified safe set, and then scale to the full 13D problem. We compare against MPC and Terminal MPC baselines, demonstrating competitive or superior docking success at lower online computational cost.
\end{enumerate}

We emphasize that the base DeepReach architecture~\cite{bansal2021deepreach} and the classical BRT formulation~\cite{mitchell2005} are prior work. Our contributions lie in the reach-avoid extension, the spacecraft docking application, and the integrated controller suite.

%%%%%%%%%%%%%%%%%%%%%%%%%%%%%%%%%%%%%%%%%%%%%%%%%%%%%%%%%%%%%%%%%%%%%%%%%%%%%%%%
\section{PROBLEM STATEMENT}

We consider a chaser spacecraft performing autonomous docking with a stationary target in a circular orbit at altitude $h = 400$\,km. The relative motion is described in the target-centered Local Vertical Local Horizontal (LVLH) frame, where $+x$ points radially outward, $+y$ is along-track, and $+z$ completes the right-hand system.

\subsection{System Dynamics}

\textit{6D Planar Model.}
For validation, we first study a planar (in-plane) docking model with state
\begin{equation}\label{eq:state6d}
x = [p_x,\; p_y,\; v_x,\; v_y,\; \theta,\; \omega]^\top \in \mathbb{R}^6,
\end{equation}
where $(p_x, p_y)$ is the relative position, $(v_x, v_y)$ the relative velocity, $\theta$ the chaser yaw angle, and $\omega$ its angular rate. The control input $u = [u_x, u_y, u_\theta]^\top$ comprises forces $|u_{x,y}| \leq \bar{F} = 20$\,N and torque $|u_\theta| \leq \bar{\tau} = 1.5$\,N$\cdot$m. The dynamics follow the Clohessy--Wiltshire (CW) equations~\cite{clohessy1960} augmented with rotational kinematics:
\begin{equation}\label{eq:cw6d}
\dot{x} = \begin{bmatrix} v_x \\ v_y \\ 3n^2 p_x + 2nv_y + u_x/m_c \\ -2nv_x + u_y/m_c \\ \omega \\ u_\theta / J_c \end{bmatrix},
\end{equation}
where $n = \sqrt{\mu/r^3} \approx 1.082 \times 10^{-3}$\,rad/s is the mean motion, $m_c = 200$\,kg the chaser mass, and $J_c = \tfrac{1}{12}m_c(w_c^2 + h_c^2) \approx 33.33$\,kg$\cdot$m$^2$ the planar moment of inertia for a $1\times1\times1$\,m cube.

\textit{13D Full Model.}
The full three-dimensional model has state
\begin{equation}\label{eq:state13d}
x = [p_x, p_y, p_z, v_x, v_y, v_z, q_0, q_1, q_2, q_3, \omega_x, \omega_y, \omega_z]^\top \!\in\! \mathbb{R}^{13}\!,
\end{equation}
where $(p_x, p_y, p_z)$ and $(v_x, v_y, v_z)$ are position and velocity in LVLH, $\mathbf{q} = [q_0, q_1, q_2, q_3]^\top$ is the scalar-first unit quaternion from LVLH to body frame, and $(\omega_x, \omega_y, \omega_z)$ are body-frame angular rates relative to LVLH. The control $u = [F_x, F_y, F_z, \tau_x, \tau_y, \tau_z]^\top$ comprises body-frame forces $|F_i| \leq 20$\,N and torques $|\tau_i| \leq 1.5$\,N$\cdot$m. The translational dynamics follow the three-dimensional Hill--Clohessy--Wiltshire (HCW) equations with body-to-LVLH force transformation, and the rotational dynamics follow Euler's equation with quaternion kinematics. The full 13D equations are given in the Appendix.

Table~\ref{tab:params} summarizes the key system parameters.

\begin{table}[t]
\caption{System Parameters}
\label{tab:params}
\centering
\begin{tabular}{lll}
\toprule
Parameter & Symbol & Value \\
\midrule
Gravitational parameter & $\mu$ & $3.986\times10^{14}$\,m$^3$/s$^2$ \\
Orbital altitude & $h$ & 400\,km \\
Mean motion & $n$ & $1.082\times10^{-3}$\,rad/s \\
Chaser mass & $m_c$ & 200\,kg \\
Chaser dimensions & $w_c\!\times\!h_c\!\times\!d_c$ & $1\!\times\!1\!\times\!1$\,m \\
Max force/axis & $\bar{F}$ & 20\,N \\
Max torque/axis & $\bar{\tau}$ & 1.5\,N$\cdot$m \\
Inertia (each axis) & $I_{xx}\!=\!I_{yy}\!=\!I_{zz}$ & 33.33\,kg$\cdot$m$^2$ \\
\bottomrule
\end{tabular}
\end{table}

\subsection{Goal and Failure Sets}

\textit{Target Geometry.}
The target body occupies $y \in [0, h_t]$ with $|x| \leq w_t/2$ (and $|z| \leq d_t/2$ in 3D), where $w_t = 6$\,m, $h_t = 3$\,m, $d_t = 3$\,m. A docking post protrudes from the target's $-y$ face at $y \in [-\ell_p, 0]$ with half-widths $0.6$\,m, where $\ell_p = 0.2$\,m is the post length.

\textit{Goal Set.}
We define the reach function $\ell(x)$ as a shaped signed-distance function, negative inside the goal and positive outside. Docking is achieved when the chaser simultaneously satisfies position, velocity, attitude, and angular rate tolerances within a goal band located just below the docking post tip. The 13D tolerances follow the IDSS~\cite{nasaidss}: lateral position $\varepsilon_p = 0.10$\,m, lateral velocity $\varepsilon_{v,\mathrm{lat}} = 0.02$\,m/s, axial approach velocity $v_y \in [0.03, 0.10]$\,m/s, attitude error $\varepsilon_q = 0.151$\,rad ($\approx 8.7^\circ$), and angular rate $\varepsilon_\omega = 0.00698$\,rad/s ($\approx 0.4^\circ$/s). The 6D tolerances are $\varepsilon_p = 0.1$\,m, $\varepsilon_v = 0.1$\,m/s, $\varepsilon_\theta = 0.04$\,rad, and $\varepsilon_\omega = 0.05$\,rad/s. The goal quaternion corresponds to a $90^\circ$ yaw so the chaser's $-y$ body axis faces the docking port.

The reach function combines shaped component distances:
\begin{equation}\label{eq:reach_fn}
\ell(x) = c_\ell \cdot \max(d_{\mathrm{pos}},\; d_{\mathrm{vel}},\; d_{\mathrm{att}},\; d_{\mathrm{rate}}),
\end{equation}
where each component $d_i$ is the signed distance to its tolerance, piecewise-shaped with linear scaling inside the goal (ensuring gradient visibility) and $\tanh$ saturation outside. The weight $c_\ell$ normalizes the reach function (1.2 for 6D, 1.0 for 13D). The inner scales are empirically tuned so that each component reaches approximately equal depth ($\approx -3$) at the goal center, giving the neural network uniform gradient visibility across all state dimensions.

\textit{Failure Set.}
The avoid function $g(x)$ is a signed-distance function, negative inside the obstacle (unsafe) and positive outside (safe). The obstacle is the union of the target body and docking post, each inflated by the chaser's collision buffer $r_c$:
\begin{equation}\label{eq:avoid_fn}
g(x) = \min\!\big(s_{\mathrm{body}}(x),\; s_{\mathrm{post}}(x)\big),
\end{equation}
where $s_{\mathrm{body}}$ and $s_{\mathrm{post}}$ are signed distances to the inflated rectangles (6D) or prisms (13D). In 6D, $r_c = \sqrt{2}/2 \approx 0.707$\,m (2D diagonal half-extent). In 13D, the buffer in the $y$-direction is orientation-dependent:
\begin{equation}\label{eq:cb_y}
r_{c,y}(\mathbf{q}) = \tfrac{1}{2}(|R_{01}| + |R_{11}| + |R_{21}|),
\end{equation}
where $R(\mathbf{q})$ is the LVLH-to-body rotation matrix. This support function ranges from 0.5\,m (aligned) to $\sqrt{3}/2 \approx 0.866$\,m (worst case), allowing the goal band to sit closer to the post when the chaser is properly oriented for docking. The failure region uses asymmetric shaping (e.g., $5\times$ scaling inside, $\tanh$ saturation outside for 13D) and is weighted by $w_g = 0.5$ to balance against the reach function.

\textit{Boundary Function.}
The combined terminal condition for the reach-avoid problem is
\begin{equation}\label{eq:boundary}
b(x) = \max\!\big(\ell(x),\; -g(x)\big).
\end{equation}
A state satisfies $b(x) \leq 0$ if and only if it is inside the goal ($\ell \leq 0$) \emph{and} outside the obstacle ($g > 0$).

   \begin{figure}[t]
      \centering
      \framebox{\parbox{3in}{Diagram of chaser--target docking geometry in the LVLH frame, showing the target body (rectangle/prism), docking post, goal band (green), failure set boundary (red), and collision buffer. Coordinate axes labeled. Both 2D (6D model) and 3D (13D model) views shown.}}
      \caption{Chaser--target docking geometry. The goal set (green) is a narrow band below the docking post tip. The failure set (red) is the target body and post, inflated by the orientation-dependent collision buffer. The chaser must navigate to the goal while avoiding the failure set.}
      \label{fig:geometry}
   \end{figure}

\subsection{Reach-Avoid Problem}

Given the dynamics $\dot{x} = f(x, u)$, goal set $\mathcal{G} = \{x : \ell(x) \leq 0\}$, and failure set $\mathcal{F} = \{x : g(x) \leq 0\}$, the \emph{Backward Reach-Avoid Tube} (BRAT) at time horizon $t$ is
\begin{equation}\label{eq:brat}
\mathcal{R}(t) = \left\{ x_0 \;\middle|\;
\begin{array}{l}
\exists\, u(\cdot) \in \mathcal{U}^{[0,t]},\;\exists\, \tau_1 \in [0,t]: \\
\xi_{x_0}^u(\tau_1) \in \mathcal{G},\;\; \xi_{x_0}^u(\tau_2) \notin \mathcal{F}\;\;\forall\, \tau_2 \in [0, \tau_1]
\end{array}
\right\},
\end{equation}
where $\xi_{x_0}^u(\cdot)$ denotes the trajectory from $x_0$ under control $u(\cdot)$. The BRAT collects all initial states from which the chaser can reach the docking goal while avoiding the target body within time horizon $t$.

%%%%%%%%%%%%%%%%%%%%%%%%%%%%%%%%%%%%%%%%%%%%%%%%%%%%%%%%%%%%%%%%%%%%%%%%%%%%%%%%
\section{PRELIMINARIES}

\subsection{HJ Value Function for BRATs}

We encode the BRAT via a value function $V: \mathbb{R}^n \times [0, T] \to \mathbb{R}$, parameterized by backward time $t \in [0, T]$ where $t = 0$ is the terminal time and $t = T$ is the maximum planning horizon. The zero sub-level set of $V$ recovers the BRAT: $\mathcal{R}(t) = \{x : V(x, t) \leq 0\}$.

The reach-avoid cost functional for a trajectory $\xi$ under control $u$ is
\begin{equation}\label{eq:cost}
J(x, u(\cdot), t) = \min_{\tau_1 \in [0,t]} \max\!\Big\{\ell\!\big(\xi(\tau_1)\big),\; \max_{\tau_2 \in [0,\tau_1]} \big\{-g\!\big(\xi(\tau_2)\big)\big\}\Big\}.
\end{equation}
The value function is the optimal cost under the best control:
\begin{equation}\label{eq:valfun}
V(x, t) = \min_{u(\cdot)} J(x, u(\cdot), t).
\end{equation}

The value function satisfies the reach-avoid variational inequality (VI)~\cite{fisac2015, margellos2011}:
\begin{equation}\label{eq:hjvi}
\min\!\Big\{\max\!\Big\{\frac{\partial V}{\partial t} - H(x, \nabla_x V),\; V - \ell(x)\Big\},\; V + g(x)\Big\} = 0,
\end{equation}
for $t \in (0, T]$, with terminal condition $V(x, 0) = b(x)$, where the Hamiltonian is
\begin{equation}\label{eq:hamiltonian}
H(x, p) = \min_{u \in \mathcal{U}} \; p^\top f(x, u).
\end{equation}
The inner $\max$ ensures either the PDE evolution holds or $V$ is capped by the reach function (goal already reached). The outer $\min$ ensures $V \geq -g$ (obstacle dominance). We note that this formulation does not include disturbances; the robust extension is deferred to future work.

\textit{Optimal Control.}
For the control-affine dynamics $f(x, u) = f_0(x) + Bu$, minimizing the Hamiltonian yields bang-bang control:
\begin{equation}\label{eq:bangbang}
u_i^* = -\bar{u}_i \cdot \mathrm{sign}\!\big((B^\top \nabla_x V)_i\big),
\end{equation}
where $\bar{u}_i$ is the $i$-th control bound. For the 13D system, body-frame force allocation requires rotation: the gradient $\nabla_v V$ in LVLH is transformed via $R(\mathbf{q})\nabla_v V$ before applying the sign operator (see Appendix).

\subsection{DeepReach}

DeepReach~\cite{bansal2021deepreach} approximates $V(x,t)$ with a neural network $\phi_\theta(x,t)$ using SIREN architecture~\cite{sitzmann2020siren}---a fully connected network with sinusoidal activations $\sigma(z) = \sin(\omega_0 z)$, where $\omega_0 = 30$. The PDE residual of~\eqref{eq:hjvi} serves as a self-supervised loss, with gradients $\partial V/\partial t$ and $\nabla_x V$ computed via automatic differentiation through the network. We extend this framework from BRT (avoid-only) to the full reach-avoid VI.

\subsection{Model Predictive Control}

We employ MPC both as a training-time label generator and as a runtime baseline. The MPC solves a finite-horizon trajectory optimization via perturbation-based shooting: at each step, $N_s$ control perturbations are sampled from a Gaussian centered on the current best trajectory, rolled out through the dynamics, and evaluated using the reach-avoid cost~\eqref{eq:cost}. The lowest-cost trajectory is retained and iteratively refined over $N_r$ rounds.

\begin{remark}
The safety characterization provided by the learned value function holds under the assumption that the neural network has converged to a solution of the VI~\eqref{eq:hjvi}. Verifying this convergence for neural approximations of high-dimensional value functions remains an open problem. We assess convergence empirically via training loss plateaus and Monte Carlo validation against grid-based ground truth in 6D.
\end{remark}

%%%%%%%%%%%%%%%%%%%%%%%%%%%%%%%%%%%%%%%%%%%%%%%%%%%%%%%%%%%%%%%%%%%%%%%%%%%%%%%%
\section{METHOD}

\subsection{Exact-Initial-Condition Parameterization}

A critical challenge in training neural value functions for reach-avoid problems is enforcing the terminal condition $V(x, 0) = b(x)$. Standard DeepReach uses a Dirichlet loss term that penalizes deviations at $t = 0$, but this boundary is never exactly satisfied. We introduce an \emph{exact} parameterization:
\begin{equation}\label{eq:exact}
V(x, t) = \phi_\theta(x, t) \cdot t \cdot \frac{\sigma_v}{\sigma_n} + b(x),
\end{equation}
where $\phi_\theta$ is the SIREN network output, $\sigma_v/\sigma_n$ is a normalization ratio (empirically tuned to $1/0.02 = 50$ for 6D and $1/0.05 = 20$ for 13D), and $b(x)$ is the boundary function~\eqref{eq:boundary}. The multiplication by $t$ ensures $V(x, 0) = b(x)$ \emph{exactly} regardless of the network weights, eliminating the Dirichlet loss entirely. This is analogous to the ``exact boundary'' trick in physics-informed neural networks~\cite{lagaris1998} but tailored here to the reach-avoid boundary structure.

The gradients under this parameterization are
\begin{align}\label{eq:exact_grads}
\frac{\partial V}{\partial t} &= \frac{\sigma_v}{\sigma_n}\!\left(t \frac{\partial \phi_\theta}{\partial t} + \phi_\theta\right)\!,\\
\nabla_x V &= \frac{\sigma_v}{\sigma_n} \cdot t \cdot \nabla_x \phi_\theta + \nabla_x b(x). \nonumber
\end{align}

\subsection{Training Pipeline}

\textit{Network Architecture.}
We use a SIREN MLP with 3 hidden layers of 512 units each. The input dimension is 8 for 6D (time + 6 states, with angle $\theta$ expanded to $(\sin\theta, \cos\theta)$ for periodicity) and 14 for 13D (time + 13 states, with quaternions canonicalized to enforce $q_0 \geq 0$). The output is a single scalar. Training uses Adam with learning rate $2 \times 10^{-5}$ and exponential decay $\gamma = 0.9997$ over 150{,}000 epochs.

\textit{Loss Function.}
The total loss combines three terms:
\begin{equation}\label{eq:loss}
\mathcal{L} = w_{\mathrm{pde}} \cdot \mathcal{L}_{\mathrm{pde}} + w_{\mathrm{dir}} \cdot \mathcal{L}_{\mathrm{dir}} + w_{\mathrm{mpc}} \cdot \mathcal{L}_{\mathrm{mpc}}.
\end{equation}
With the exact parameterization~\eqref{eq:exact}, $\mathcal{L}_{\mathrm{dir}} = 0$ by construction.

The \emph{PDE loss} penalizes violations of the reach-avoid VI~\eqref{eq:hjvi}:
\begin{equation}\label{eq:pde_loss}
\mathcal{L}_{\mathrm{pde}} = \frac{1}{N}\sum_{i=1}^{N} \left| \min\!\left\{\max\!\left\{\frac{\partial V}{\partial t} - H,\; V - \ell\right\},\; V + g\right\}\right|,
\end{equation}
evaluated at $N = 50{,}000$ collocation points per batch.

The \emph{MPC supervision loss} aligns the learned value with MPC-generated reach-avoid cost labels:
\begin{equation}\label{eq:mpc_loss}
\mathcal{L}_{\mathrm{mpc}} = \frac{1}{M}\sum_{j=1}^{M} |V_\theta(x_j, t_j) - V_{\mathrm{mpc}}(x_j, t_j)|,
\end{equation}
with an additional false-positive penalty: when the MPC label indicates unsafe ($V_{\mathrm{mpc}} > 0$) but the network predicts safe ($V_\theta \leq 0$), a penalty term $\lambda |V_\theta - V_{\mathrm{mpc}}| \cdot |V_\theta|$ is added, with $\lambda = 100$ to strongly discourage under-approximation of the unsafe region.

\textit{Curriculum Learning.}
Training proceeds with a time-based curriculum: the sampled time horizon grows linearly from $0$ to $T$ over the first $100{,}000$ epochs:
\begin{equation}\label{eq:curriculum}
t_{\max}(\text{epoch}) = T \cdot 1.1 \cdot \min\!\left(\frac{\text{epoch}}{100{,}000},\; 1\right)\!,
\end{equation}
where $T = 15$\,s for 6D and $T = 10$\,s for 13D. The factor of $1.1$ provides slight overshoot for numerical stability. Training can be paused at intermediate horizons to allow value function convergence before extending further, assessed by comparing the learned $V$ against MPC ground truth on held-out states.

\textit{MPC Label Generation.}
MPC labels are generated by rolling out $N_s = 100$ Gaussian-perturbed control trajectories from each of $1{,}000$ initial conditions, retaining the lowest-cost trajectory and computing the reach-avoid cost-to-go at each timestep. Labels are regenerated periodically (every 10{,}000--20{,}000 epochs) to track the evolving value function. A multi-tier sampling strategy concentrates training points near the goal region: 10--15\% exact goal states with small noise, 30\% Gaussian around the goal at $2\times$ tolerance width, 20--25\% boundary-focused samples at $0.8$--$1.2\times$ tolerance, and 30--40\% broader uniform samples.

\subsection{Online Controller Suite}

We implement four controllers for comparison. All operate at $\Delta t = 0.1$\,s with Euler integration.

\textit{BRAT Controller} (our primary contribution).
A two-phase bang-bang controller driven by the learned value function gradient.
\emph{Phase~1 (Convergence)}: While $V(x, T) > 0$ (state outside the BRAT), query the value function at fixed horizon $t = T$ and apply bang-bang control~\eqref{eq:bangbang}. If the gradient magnitude on control-coupled dimensions falls below a threshold ($0.01$), a virtual PD gradient is blended in as a fallback.
\emph{Phase~2 (Precision)}: Once $V(x, T) \leq 0$ (state enters the BRAT), perform a \emph{minimum-time search}: find the smallest $t^*$ such that $V(x, t^*) \leq 0$ and apply bang-bang control at this tighter time slice. This yields more precise terminal guidance, tightening the time horizon as the chaser approaches docking.
The transition from Phase~1 to Phase~2 is one-way.

\textit{Terminal MPC Controller.}
Short-horizon (2\,s) gradient-based MPC using the learned value function as a differentiable terminal cost. Combines the computational advantages of short-horizon planning with the long-horizon reach-avoid information encoded in $V$. Includes a stagnation-escape mechanism: if position improvement falls below $0.1$\,m over a 5\,s window, the controller escalates through exploration and BRT-fallback modes.

\textit{MPC Baseline.}
Receding-horizon MPC without any learned value function. Uses multi-start gradient shooting (6D) or perturbation-based optimization (13D) with a reach-avoid cost function over a 20\,s planning horizon.

\textit{Grid-Based Controller} (6D only).
Numerical ground truth via the \texttt{hj\_reachability} library~\cite{hjreachability}. The 6D problem is decomposed into a 4D translational subsystem ($51\!\times\!51\!\times\!31\!\times\!31$ grid) and a 2D rotational subsystem ($361\!\times\!141$ grid). The combined value function is $V_{\mathrm{6D}} = \max(V_{\mathrm{4D}}, V_{\mathrm{2D}})$, solved backward over 30\,s. Control is extracted via finite-difference gradients. This decomposition is feasible because the CW translational and single-axis rotational dynamics are decoupled in the planar model.

\textit{BRT Safety Filter.}
An optional safety override layer using a separately trained avoid-only BRT. When the avoid-only value function drops below a margin ($0.1$ in Phase~1, $0.02$ in Phase~2), the controller is overridden with bang-bang control that maximizes the avoid value, driving the chaser away from the obstacle.

   \begin{figure}[t]
      \centering
      \framebox{\parbox{3in}{Block diagram of the two-phase BRAT controller. Phase~1: query $V(x, T)$, compute $\nabla_x V$, apply bang-bang control with PD fallback. Phase~2: minimum-time search for $t^*$, query $V(x, t^*)$, apply tighter bang-bang control. Transition triggered when $V(x, T) \leq 0$.}}
      \caption{BRAT controller architecture. The controller transitions from convergence (Phase~1) to precision (Phase~2) when the state enters the BRAT zero sub-level set.}
      \label{fig:controller}
   \end{figure}

%%%%%%%%%%%%%%%%%%%%%%%%%%%%%%%%%%%%%%%%%%%%%%%%%%%%%%%%%%%%%%%%%%%%%%%%%%%%%%%%
\section{RESULTS}

\subsection{Experimental Setup}

We evaluate all controllers via closed-loop Monte Carlo simulations. Initial conditions are sampled uniformly from a reduced state range (e.g., $|p_{x,y}| \leq 13$\,m, $|v_{x,y}| \leq 0.75$\,m/s for 6D) and filtered to exclude states inside the obstacle or already at the goal. Each simulation runs until docking success, collision, or a timeout of [TODO]\,s.

We report: \emph{docking rate} (successful dockings / total), \emph{collision rate}, \emph{timeout rate}, \emph{mean docking time}, \emph{mean control effort} $\sum \|u\| \Delta t$ over successful trials, and \emph{mean wall-clock time} per decision step. All experiments were performed using [TODO: specify computing environment].

\subsection{6D Validation Against Grid-Based Ground Truth}

   \begin{figure}[t]
      \centering
      \framebox{\parbox{3in}{2D contour plots comparing the learned BRAT value function (DeepReach) and the grid-based value function at three state slices: (a) position $(p_x, p_y)$ at $v_x\!=\!v_y\!=\!0$, $\theta\!=\!\pi/2$, $\omega\!=\!0$; (b) velocity $(v_x, v_y)$ at $p_x\!=\!0$, $p_y\!\approx\!{-1.2}$\,m; (c) attitude $(\theta, \omega)$ at nominal translational state. Zero-level contours of both value functions overlaid. Goal set boundary in green, failure set boundary in red.}}
      \caption{Value function comparison between the learned BRAT and grid-based ground truth in 6D. The learned zero-level set (solid) closely approximates the grid-based solution (dashed) across position, velocity, and attitude slices.}
      \label{fig:value_comparison}
   \end{figure}

We validate the learned 6D value function against the decomposed grid-based solution. Monte Carlo volume estimation ($N = 500{,}000$--$2{,}000{,}000$ samples) over the evaluation region $p_{x,y} \in [-5, 5]$\,m, $v_{x,y} \in [-1.5, 1.5]$\,m/s, $\theta \in [-\pi, \pi]$, $\omega \in [-1, 1]$\,rad/s yields:

\begin{table}[t]
\caption{6D BRAT Volume Overlap with Grid-Based Ground Truth}
\label{tab:volume}
\centering
\begin{tabular}{lc}
\toprule
Metric & Value \\
\midrule
Grid BRAT volume (fraction of eval.\ region) & [TODO] \\
Learned BRAT volume & [TODO] \\
Overlap volume & [TODO] \\
Jaccard index ($|$overlap$|/|$union$|$) & [TODO] \\
Gradient sign agreement (control-coupled dims.) & [TODO] \\
Cosine similarity ($\nabla V_{\mathrm{grid}},\nabla V_{\mathrm{DR}}$) & [TODO] \\
\bottomrule
\end{tabular}
\end{table}

Since bang-bang control depends only on the sign of the value function gradient, we also compare gradient sign agreement on the control-coupled dimensions ($v_x$, $v_y$, $\omega$). A high sign agreement rate indicates that the learned and grid-based controllers produce identical control actions at most states.

\subsection{Controller Comparison}

\textit{6D Results.}
Table~\ref{tab:6d_results} reports performance across four controllers on [TODO] shared initial conditions.

\begin{table}[t]
\caption{6D Controller Comparison ($N = $ [TODO] rollouts)}
\label{tab:6d_results}
\centering
\begin{tabular}{lcccc}
\toprule
Metric & BRAT & Grid & MPC & T-MPC \\
\midrule
Docking rate (\%) & [TODO] & [TODO] & [TODO] & [TODO] \\
Collision rate (\%) & [TODO] & [TODO] & [TODO] & [TODO] \\
Timeout rate (\%) & [TODO] & [TODO] & [TODO] & [TODO] \\
Mean dock time (s) & [TODO] & [TODO] & [TODO] & [TODO] \\
Mean effort (N$\cdot$s) & [TODO] & [TODO] & [TODO] & [TODO] \\
Mean step time (ms) & [TODO] & [TODO] & [TODO] & [TODO] \\
\bottomrule
\end{tabular}
\end{table}

   \begin{figure}[t]
      \centering
      \framebox{\parbox{3in}{Position-space trajectory overlays from Monte Carlo rollouts in the 6D $(p_x, p_y)$ plane. Multiple initial conditions shown for each controller (BRAT, Grid-Based, MPC, Terminal MPC) with color-coded outcomes: green for successful docking, red for collision, gray for timeout. Target body and docking post geometry shown.}}
      \caption{Representative 6D trajectories in position space. The BRAT controller (top left) produces trajectories qualitatively similar to the grid-based ground truth (top right), both navigating around the obstacle to the docking port.}
      \label{fig:6d_trajectories}
   \end{figure}

\textit{13D Results.}
Table~\ref{tab:13d_results} reports performance for the three controllers applicable to the 13D problem. The grid-based controller is excluded because the 13D state space is computationally intractable for grid-based methods---a $51$-point grid per dimension would require $51^{13} \approx 10^{22}$ grid points.

\begin{table}[t]
\caption{13D Controller Comparison ($N = $ [TODO] rollouts)}
\label{tab:13d_results}
\centering
\begin{tabular}{lccc}
\toprule
Metric & BRAT & MPC & T-MPC \\
\midrule
Docking rate (\%) & [TODO] & [TODO] & [TODO] \\
Collision rate (\%) & [TODO] & [TODO] & [TODO] \\
Timeout rate (\%) & [TODO] & [TODO] & [TODO] \\
Mean dock time (s) & [TODO] & [TODO] & [TODO] \\
Mean effort (N$\cdot$s) & [TODO] & [TODO] & [TODO] \\
Mean step time (ms) & [TODO] & [TODO] & [TODO] \\
\bottomrule
\end{tabular}
\end{table}

   \begin{figure}[t]
      \centering
      \framebox{\parbox{3in}{13D Monte Carlo results visualization. Options: (a) 3D position-space trajectory bundle $(p_x, p_y, p_z)$ with chaser/target geometry; (b) bar chart comparing docking rate, collision rate, and timeout rate across the three controllers; (c) docking time histograms for each controller.}}
      \caption{13D controller performance. [TODO: describe key findings.]}
      \label{fig:13d_results}
   \end{figure}

\subsection{Computational Performance}

The BRAT controller requires a single forward pass through the SIREN network and one backward pass for the gradient at each control step, with no online optimization. The Terminal MPC controller adds a short-horizon gradient optimization on top. Table~\ref{tab:timing} compares per-step wall-clock times.

\begin{table}[t]
\caption{Per-Step Computation Time}
\label{tab:timing}
\centering
\begin{tabular}{lcc}
\toprule
Controller & 6D (ms/step) & 13D (ms/step) \\
\midrule
BRAT & [TODO] & [TODO] \\
Terminal MPC & [TODO] & [TODO] \\
MPC & [TODO] & [TODO] \\
Grid-Based & [TODO] & N/A \\
\bottomrule
\end{tabular}
\end{table}

\addtolength{\textheight}{-5cm}

%%%%%%%%%%%%%%%%%%%%%%%%%%%%%%%%%%%%%%%%%%%%%%%%%%%%%%%%%%%%%%%%%%%%%%%%%%%%%%%%
\section{CONCLUSION}

We presented a learning-based BRAT framework for autonomous spacecraft docking in 13 dimensions, combining neural value function approximation with MPC-generated curriculum labels and an exact-initial-condition parameterization tailored to reach-avoid semantics.

\textit{Limitations.}
First, the neural value function lacks formal convergence guarantees to the true HJ solution; our safety assessment is empirical, validated via Monte Carlo rollouts and, in 6D, comparison against grid-based ground truth.
Second, the formulation is disturbance-free: we do not model thrust uncertainty, navigation errors, or unmodeled dynamics. The grid-based framework supports bounded disturbances via the differential game formulation, but this extension remains to be incorporated into the neural approach.
Third, the target is assumed stationary in its LVLH frame; docking with tumbling or actively maneuvering targets is not addressed.

\textit{Future Work.}
To address the convergence limitation, we plan to investigate neural network verification techniques adapted to high-dimensional value functions, complementing the empirical validation with formal certificates.
Robustness to disturbances can be incorporated by extending the Hamiltonian to include a maximizing adversary, following the standard differential game formulation~\cite{mitchell2005}, with the neural training pipeline adapted accordingly.
Extension to tumbling targets requires a time-varying goal and failure set, addressable within the HJ framework by coupling target attitude dynamics into the state~\cite{mote2023}.

%%%%%%%%%%%%%%%%%%%%%%%%%%%%%%%%%%%%%%%%%%%%%%%%%%%%%%%%%%%%%%%%%%%%%%%%%%%%%%%%
\section*{APPENDIX: FULL 13D DYNAMICS}

The 13D dynamics consist of three-dimensional HCW translational equations, quaternion kinematics, and Euler's rotational equation.

\textit{Translation.}
Forces applied in body frame are transformed to LVLH acceleration via the rotation matrix $R(\mathbf{q})$:
\begin{equation}\label{eq:13d_trans}
\mathbf{a}_L = \frac{1}{m_c} R(\mathbf{q})^\top \mathbf{F}_b,
\end{equation}
where $\mathbf{F}_b = [F_x, F_y, F_z]^\top$ is the body-frame force vector and $R(\mathbf{q})$ is the LVLH-to-body rotation from the scalar-first quaternion. The HCW accelerations are
\begin{align}\label{eq:hcw3d}
\dot{v}_x &= 2nv_y + 3n^2 p_x + a_{L,x}, \nonumber\\
\dot{v}_y &= -2nv_x + a_{L,y}, \\
\dot{v}_z &= -n^2 p_z + a_{L,z}. \nonumber
\end{align}

\textit{Quaternion Kinematics.}
The quaternion evolves according to
\begin{equation}\label{eq:qdot}
\dot{\mathbf{q}} = \tfrac{1}{2}\,\Omega(\bm{\omega})\,\mathbf{q},
\end{equation}
where $\bm{\omega} = [\omega_x, \omega_y, \omega_z]^\top$ is the body-frame angular velocity relative to LVLH and the $4\times4$ matrix is
\begin{equation}\label{eq:omega_mat}
\Omega(\bm{\omega}) = \begin{bmatrix}
0 & -\omega_x & -\omega_y & -\omega_z \\
\omega_x & 0 & \omega_z & -\omega_y \\
\omega_y & -\omega_z & 0 & \omega_x \\
\omega_z & \omega_y & -\omega_x & 0
\end{bmatrix}\!.
\end{equation}
To resolve the $\mathbf{q} \leftrightarrow -\mathbf{q}$ ambiguity, we enforce $q_0 \geq 0$ by canonicalization before network evaluation.

\textit{Rotational Dynamics.}
Euler's equation for the angular velocity is
\begin{equation}\label{eq:euler}
\dot{\bm{\omega}} = I^{-1}\!\left(\bm{\tau}_b - \bm{\omega} \times I\bm{\omega}\right)\!,
\end{equation}
where $I = \mathrm{diag}(I_{xx}, I_{yy}, I_{zz})$ with $I_{xx} = I_{yy} = I_{zz} \approx 33.33$\,kg$\cdot$m$^2$ and $\bm{\tau}_b = [\tau_x, \tau_y, \tau_z]^\top$ is the body-frame torque.

\textit{13D Optimal Control.}
The force allocation accounts for orientation:
\begin{align}\label{eq:13d_control}
\mathbf{c}_F &= \tfrac{1}{m_c}\,R(\mathbf{q})\,\nabla_{\mathbf{v}} V, & F_i^* &= -\bar{F}\cdot\mathrm{sign}(c_{F,i}), \nonumber\\
\mathbf{c}_\tau &= I^{-\!\top}\!\nabla_{\bm{\omega}} V, & \tau_i^* &= -\bar{\tau}\cdot\mathrm{sign}(c_{\tau,i}).
\end{align}

\section*{ACKNOWLEDGMENT}

[TODO: Add acknowledgments.]

%%%%%%%%%%%%%%%%%%%%%%%%%%%%%%%%%%%%%%%%%%%%%%%%%%%%%%%%%%%%%%%%%%%%%%%%%%%%%%%%

\begin{thebibliography}{99}

\bibitem{mitchell2005}
I.~M. Mitchell, A.~M. Bayen, and C.~J. Tomlin, ``A time-dependent Hamilton--Jacobi formulation of reachable sets for continuous dynamic games,'' \emph{IEEE Trans. Autom. Control}, vol.~50, no.~7, pp.~947--957, Jul.~2005.

\bibitem{bansal2017hjsurvey}
S.~Bansal, M.~Chen, S.~Herbert, and C.~J. Tomlin, ``Hamilton--Jacobi reachability: A brief overview and recent advances,'' in \emph{Proc. IEEE Conf. Decision Control (CDC)}, 2017, pp.~2242--2253.

\bibitem{bansal2021deepreach}
S.~Bansal and C.~J. Tomlin, ``DeepReach: A deep learning approach to high-dimensional reachability,'' in \emph{Proc. IEEE Int. Conf. Robot. Autom. (ICRA)}, 2021, pp.~1817--1824.

\bibitem{fisac2015}
J.~F. Fisac, M.~Chen, C.~J. Tomlin, and S.~S. Sastry, ``Reach-avoid problems with time-varying dynamics, targets and constraints,'' in \emph{Proc. 18th Int. Conf. Hybrid Syst.: Comput. Control (HSCC)}, 2015, pp.~11--20.

\bibitem{clohessy1960}
W.~H. Clohessy and R.~S. Wiltshire, ``Terminal guidance system for satellite rendezvous,'' \emph{J. Aerospace Sci.}, vol.~27, no.~9, pp.~653--658, Sep.~1960.

\bibitem{sitzmann2020siren}
V.~Sitzmann, J.~N.~P. Martel, A.~W. Bergman, D.~B. Lindell, and G.~Wetzstein, ``Implicit neural representations with periodic activation functions,'' in \emph{Advances Neural Inf. Process. Syst. (NeurIPS)}, vol.~33, 2020, pp.~7462--7473.

\bibitem{chen2018decomposition}
M.~Chen, S.~L. Herbert, M.~S. Vashishtha, S.~Bansal, and C.~J. Tomlin, ``Decomposition of reachable sets and tubes for a class of nonlinear systems,'' \emph{IEEE Trans. Autom. Control}, vol.~63, no.~11, pp.~3675--3688, Nov.~2018.

\bibitem{fehse2003}
W.~Fehse, \emph{Automated Rendezvous and Docking of Spacecraft}.\hskip 1em plus 0.5em minus 0.4em Cambridge, U.K.: Cambridge Univ. Press, 2003.

\bibitem{nasaidss}
NASA, ``International Docking System Standard (IDSS) Interface Definition Document,'' NASA, Tech. Rep. Rev.~F, 2022.

\bibitem{hjreachability}
Stanford ASL, ``hj\_reachability: Hamilton--Jacobi reachability in JAX,'' 2021. [Online]. Available: https://github.com/StanfordASL/hj\_reachability

\bibitem{malyuta2022}
D.~Malyuta, T.~P. Reynolds, M.~Szmuk, T.~Lew, R.~Bonalli, M.~Pavone, and B.~A\c{c}{\i}kme\c{s}e, ``Convex optimization for trajectory generation: A tutorial on generating dynamically feasible trajectories reliably and efficiently,'' \emph{IEEE Control Syst. Mag.}, vol.~42, no.~5, pp.~40--113, Oct.~2022.

\bibitem{ames2019cbf}
A.~D. Ames, S.~Coogan, M.~Egerstedt, G.~Notomista, K.~Sreenath, and P.~Tabuada, ``Control barrier functions: Theory and applications,'' in \emph{Proc. Eur. Control Conf. (ECC)}, 2019, pp.~3420--3431.

\bibitem{wabersich2021}
K.~P. Wabersich and M.~N. Zeilinger, ``A predictive safety filter for learning-based control of constrained nonlinear dynamical systems,'' \emph{Automatica}, vol.~129, 109597, Jul.~2021.

\bibitem{jewison2016}
C.~Jewison and R.~S. Erwin, ``A spacecraft benchmark problem for hybrid control and estimation,'' in \emph{Proc. IEEE Conf. Decision Control (CDC)}, 2016, pp.~3300--3305.

\bibitem{ventura2017}
J.~Ventura, M.~Romano, and U.~Walter, ``A real-time computational framework for guidance, navigation, and control of spacecraft proximity maneuvers,'' \emph{Acta Astronautica}, vol.~132, pp.~30--44, Mar.~2017.

\bibitem{bui2022}
T.~Bui, I.~M. Mitchell, and D.~J. Scheeres, ``Computing reach-avoid sets for space vehicle docking under continuous thrust,'' in \emph{Proc. IEEE Aerosp. Conf.}, 2022. [TODO: verify]

\bibitem{mote2023}
M.~Mote, M.~Chen, K.~Dunlap, and K.~Kim, ``Applied reachability analysis for spacecraft rendezvous and docking with a tumbling object,'' [TODO: find venue], 2023.

\bibitem{hovell2021}
K.~Hovell and S.~Ulrich, ``Deep reinforcement learning for spacecraft proximity operations guidance,'' \emph{J. Spacecraft Rockets}, vol.~58, no.~2, pp.~254--264, 2021. [TODO: verify]

\bibitem{broida2019}
J.~Broida and R.~Linares, ``Spacecraft rendezvous guidance in cluttered environments via reinforcement learning,'' in \emph{Proc. AAS/AIAA Space Flight Mech. Meeting}, 2019. [TODO: verify]

\bibitem{eren2015}
U.~Eren, A.~Prach, B.~B. Koco\u{g}lu, S.~Akg\"{o}l, F.~Inalhan, and M.~A. Akg\"{u}l, ``Model predictive control in aerospace systems: Current state and opportunities,'' \emph{J. Guid. Control Dyn.}, vol.~40, no.~7, pp.~1541--1566, 2017. [TODO: verify]

\bibitem{margellos2011}
K.~Margellos and J.~Lygeros, ``Hamilton--Jacobi formulation for reach-avoid differential games,'' \emph{IEEE Trans. Autom. Control}, vol.~56, no.~8, pp.~1849--1861, Aug.~2011.

\bibitem{lygeros2004}
J.~Lygeros, ``On reachability and minimum cost optimal control,'' \emph{Automatica}, vol.~40, no.~6, pp.~917--927, Jun.~2004.

\bibitem{lagaris1998}
I.~E. Lagaris, A.~Likas, and D.~I. Fotiadis, ``Artificial neural networks for solving ordinary and partial differential equations,'' \emph{IEEE Trans. Neural Netw.}, vol.~9, no.~5, pp.~987--1000, Sep.~1998.

\bibitem{borquez2023}
J.~Borquez, S.~Aggarwal, and S.~L. Herbert, ``Parameter-conditioned reachable sets for updating safety assurances online,'' in \emph{Proc. IEEE Int. Conf. Robot. Autom. (ICRA)}, 2023. [TODO: verify]

\bibitem{shao2023}
Y.~M. Shao and S.~Bansal, ``Reachability-based trajectory design with neural implicit safety constraints,'' in \emph{Proc. Robot.: Sci. Syst. (RSS)}, 2023. [TODO: verify]

\bibitem{hwo2024}
National Academies, ``Pathways to Discovery in Astronomy and Astrophysics for the 2020s (Astro2020 Decadal Survey),'' 2021. [TODO: verify HWO reference]

\bibitem{pytorch}
A.~Paszke \emph{et al.}, ``PyTorch: An imperative style, high-performance deep learning library,'' in \emph{Advances Neural Inf. Process. Syst. (NeurIPS)}, vol.~32, 2019, pp.~8026--8037.

\bibitem{hsu2021}
S.~C. Hsu, J.~F. Rubies-Royo, and C.~J. Tomlin, ``Safety and liveness guarantees through reach-avoid reinforcement learning,'' in \emph{Proc. Robot.: Sci. Syst. (RSS)}, 2021. [TODO: verify]

\bibitem{leung2020}
K.~Leung, E.~Schmerling, M.~Chen, J.~Shim, and M.~Pavone, ``Robust safety guarantees for uncertain nonlinear systems via Hamilton--Jacobi reachability analysis,'' \emph{Int. J. Robot. Res.}, vol.~39, no.~13, pp.~1550--1570, 2020. [TODO: verify]

\end{thebibliography}

\end{document}
